This study asks to what extent AEMO five-minute intra-hour features improve short-term solar PV forecast skill over meteorological inputs alone in New South Wales. Five-minute AEMO observations were aggregated into complete hours and aligned with regional hourly MERRA-2 reanalysis. Exact duplicate timestamps were collapsed and conflicting duplicates removed. XGBoost, LSTM, an unpenalised CNN-BiLSTM and a physics-guided CNN-BiLSTM were compared with persistence, rolling climatology and SARIMA at 1, 3, 6 and 24-hour horizons using a chronological 2025 test set. Neural headline forecasts average five seeds.

Without hourly generation history, the two scale-normalised dispersion summaries improve all four learned procedures robustly at every horizon, by 1.26 to 11.46 percent. Once hourly history is controlled, all four procedures gain robustly at one hour, XGBoost and the physics-guided model also gain at three and six hours, and none gains robustly at 24 hours. Robust controlled improvements range from 1.63 to 2.86 percent and require agreement across block intervals, corrected Diebold-Mariano tests and an output-scale sensitivity check.

On combined inputs, LSTM achieves the best one-hour RMSE of 155.25 MW, the physics-guided model leads at three and six hours with 300.28 and 369.17 MW, and the unpenalised CNN-BiLSTM leads at 24 hours with 399.77 MW. A matched comparison attributes a robust 1.18 percent six-hour improvement to the physics-guided loss penalties, but finds no robust penalty effect at the other horizons. The weighted ensemble does not outperform the best neural model.

The evidence supports five-minute variability as a modest short-horizon supplement, not the main source of predictive skill. Conclusions are limited to one held-out year, retrospective weather inputs and utility-scale market-cleared generation.
